% Options for packages loaded elsewhere
\PassOptionsToPackage{unicode}{hyperref}
\PassOptionsToPackage{hyphens}{url}
%
\documentclass[
  10pt,
]{article}
\usepackage{amsmath,amssymb}
\usepackage{iftex}
\ifPDFTeX
  \usepackage[T1]{fontenc}
  \usepackage[utf8]{inputenc}
  \usepackage{textcomp} % provide euro and other symbols
\else % if luatex or xetex
  \usepackage{unicode-math} % this also loads fontspec
  \defaultfontfeatures{Scale=MatchLowercase}
  \defaultfontfeatures[\rmfamily]{Ligatures=TeX,Scale=1}
\fi
\usepackage{lmodern}
\ifPDFTeX\else
  % xetex/luatex font selection
  \setmainfont[]{DejaVu Serif}
  \setmonofont[]{DejaVu Sans Mono}
\fi
% Use upquote if available, for straight quotes in verbatim environments
\IfFileExists{upquote.sty}{\usepackage{upquote}}{}
\IfFileExists{microtype.sty}{% use microtype if available
  \usepackage[]{microtype}
  \UseMicrotypeSet[protrusion]{basicmath} % disable protrusion for tt fonts
}{}
\makeatletter
\@ifundefined{KOMAClassName}{% if non-KOMA class
  \IfFileExists{parskip.sty}{%
    \usepackage{parskip}
  }{% else
    \setlength{\parindent}{0pt}
    \setlength{\parskip}{6pt plus 2pt minus 1pt}}
}{% if KOMA class
  \KOMAoptions{parskip=half}}
\makeatother
\usepackage{xcolor}
\usepackage[margin=0.82in]{geometry}
\setlength{\emergencystretch}{3em} % prevent overfull lines
\providecommand{\tightlist}{%
  \setlength{\itemsep}{0pt}\setlength{\parskip}{0pt}}
\setcounter{secnumdepth}{-\maxdimen} % remove section numbering
\ifLuaTeX
  \usepackage{selnolig}  % disable illegal ligatures
\fi
\usepackage{bookmark}
\IfFileExists{xurl.sty}{\usepackage{xurl}}{} % add URL line breaks if available
\urlstyle{same}
\hypersetup{
  hidelinks,
  pdfcreator={LaTeX via pandoc}}

\author{}
\date{}

\begin{document}

{
\setcounter{tocdepth}{3}
\tableofcontents
}
\section{Edge General Position Numbers of Lucas Cubes: A
Computer-Assisted Exact
Formula}\label{edge-general-position-numbers-of-lucas-cubes-a-computer-assisted-exact-formula}

Ryutaro Yonezu\\
Independent Researcher\\
Research No.76, v1.0, 19 November 2026

\subsection{Abstract}\label{abstract}

For each integer n \textgreater= 3, the edge general position number of
the Lucas cube Lambda\_n is F\_n + F\_\{n-2\}, with F\_0 = 0 and F\_1 =
1. The proof combines an analytic lower construction, star-forest and
projection lemmas, local capacity inequalities, finite integer
certificates, and a recurrence induction. It does not claim uniqueness
of a maximum set and does not depend on the main theorem of Research
No.75. Required exact whole-graph bases are n = 3, 4, 6. The 9,908
capacity bases, including overlaps, remain in the frozen verification
lineage. The n = 6 relaxation margin of -2 is handled by the exact
optimum 11.

\subsection{Status and scope}\label{status-and-scope}

Computer-assisted preprint. Independent AI mathematical audit: PACKAGE /
DELTA / REQUIRED BASES / INDUCTION PASS; CAPACITY BASES PASS within the
previously documented independent check range; THEOREM conditional PASS.
The auditor's historical PUBLICATION HOLD recommendation is preserved as
a record. The author separately authorized publication of this
unrefereed preprint after the publication gates passed. NOT HUMAN PEER
REVIEWED. NOVELTY INCONCLUSIVE. REFERENCES PARTIAL. Copyright (c) 2026
Ryutaro Yonezu. All rights reserved. No Creative Commons or open-source
license is granted.

The 2023 paper of Klavzar and Tan established the Lucas-cube lower bound
2F\_\{n-1\} for n \textgreater= 4 and the exact value 7 at n = 5; it did
not state the all-dimension Lucas-cube formula as an explicit
conjecture. The present classification does not claim that the n = 5
value is new. See S. Klavzar and E. Tan, \emph{Edge General Position
Sets in Fibonacci and Lucas Cubes}, Bulletin of the Malaysian
Mathematical Sciences Society 46:120 (2023), DOI:
10.1007/s40840-023-01517-y.

\subsection{Editorial note on the frozen
proof}\label{editorial-note-on-the-frozen-proof}

The following four frozen mathematical notes are reproduced verbatim, in
chronological order, from the v0.1-v0.3 lineage and v0.4 dependency
addendum. Their opening PRIVATE HOLD and audit-pending labels document
the historical stages; the current publication status is the one stated
above. The source files and original verifiers are preserved
byte-for-byte in the accompanying reproducibility archive. The v0.4 note
clarifies the final proof dependencies, without changing the earlier
mathematics. The finite verifier's historical PENDING status string is
also preserved.

\newpage

\section{Frozen v0.1: lower construction and star
capacity}\label{frozen-v0.1-lower-construction-and-star-capacity}

\section{No.76 candidate v0.1 --- A lower bound and star-capacity upper
bounds for edge general position in Lucas
cubes}\label{no.76-candidate-v0.1-a-lower-bound-and-star-capacity-upper-bounds-for-edge-general-position-in-lucas-cubes}

PRIVATE HOLD. Preliminary results, not a completed all-dimension
classification. NOVELTY INCONCLUSIVE / REFERENCES PARTIAL. No
independent audit, Windows test, or human peer review performed. No.75
is unchanged; Section 2 below is reproduced as a self-contained
argument, so this result does not depend on acceptance of No.75's
Fibonacci cube classification.

\subsection{1. Model and lower bound}\label{model-and-lower-bound}

Let Lambda\_n, n\textgreater=3, be the graph of length-n binary words
with no cyclically adjacent ones, with adjacency by a single-coordinate
flip. Let F\_0=0, F\_1=1 and F\_\{r+2\}=F\_\{r+1\}+F\_r. The cardinality
of a largest edge general position set (at most two selected edges on
each geodesic) is gp\_e. Graph distance equals Hamming distance: remove
the ones absent from the target, then insert the missing target ones.
Consequently the colors (changed coordinates) on a geodesic are
distinct. The number of vertices is the Lucas number
F\_\{n-1\}+F\_\{n+1\}.

Klavzar and Tan (2023), Proposition 4.1(ii) and Theorem 4.6, give the
coordinate-class size F\_\{n-1\} and lower bound 2F\_\{n-1\} for
n\textgreater=4. Their Proposition 4.7 establishes gp\_e(Lambda\_5)=7,
showing the lower bound 6 need not be optimal. These results are prior
work.

\textbf{Theorem A (analytic lower bound).} For every n\textgreater=4,

\begin{verbatim}
gp_e(Lambda_n) >= K_n := F_n+F_{n-2} = F_{n-1}+2F_{n-2}.
\end{verbatim}

For n\textgreater=5, K\_n=2F\_\{n-1\}+F\_\{n-4\}\textgreater2F\_\{n-1\}.

Proof. Choose three consecutive cyclic coordinates, say 0,1,2. Select
every edge changing one of the three whose other two bits are both zero.
Its projection is an arm of the 3-star centered at 000. If a geodesic
contains three selected edges, it must change each of the three
coordinates once and these are the only changes to its projection. But
the resulting 3-step projected walk would use three distinct edges of a
star while changing distinct coordinates, impossible: a simple path in a
star has length at most two. For the middle coordinate, its two
neighbors are fixed zero and the other n-3 coordinates form a free path,
giving F\_\{n-1\} edges. For either outer coordinate, the chosen three
and one additional neighbor are fixed (before the flip) to zero, leaving
a free path of n-4 coordinates, giving F\_\{n-2\} edges each, including
n=4. The sum is K\_n.~For n\textgreater=5 the Fibonacci recurrence gives
the displayed strict improvement. For n=3 the same construction selects
all three edges, and K\_3=3. QED.

\subsection{2. Two star lemmas (self-contained; same arguments as in
No.75
v0.10)}\label{two-star-lemmas-self-contained-same-arguments-as-in-no.75-v0.10}

\textbf{Lemma B.} In any isometric subgraph of a hypercube, an edge
general position set S using at least three colors induces a star forest
on its incident vertices (only selected edges are included).

Proof. A simple 3-edge path in S is forbidden if its colors are
distinct. Otherwise its colors are i,j,i. The four vertices realize all
four assignments on i,j and agree elsewhere. The oriented length-2
subpaths have all four vertices as possible ending vertices. Choose a
selected edge f of a third color k. Choose such a subpath ending at z
whose i,j bits agree with the fixed bits of f.~Choose endpoint y of f
with its k bit equal to that of z. A shortest connector z to y changes
none of i,j,k. Concatenating the length-2 subpath, connector, and f
gives a walk with distinct changed coordinates throughout, hence a
geodesic containing three selected edges, contradiction. The selected
graph is bipartite and has no simple path of length 3. Every connected
nonempty component must then be a star: two vertices of degree at least
2 would produce adjacent vertices of degree at least 2 (along a shortest
connecting path), and these produce a simple 3-edge path in a bipartite
graph. QED.

\textbf{Lemma C.} Suppose S contains a star centered at v with arms of
colors I, where \textbar I\textbar=d\textgreater=3. Every selected edge
f of color k has fixed bits on I excluding k either all equal to v
there, or all opposite to v there.

Proof. Otherwise choose i,j in I excluding k for which exactly one bit
disagrees with v. Orient the path v\textsuperscript{i--v--v}j so its
ending vertex agrees with f on i,j. Choose the endpoint of f that agrees
with this ending vertex on k. A shortest connector avoids i,j,k, and the
concatenated walk has no repeated changed coordinate, giving a forbidden
geodesic. QED.

\subsection{3. Star-capacity upper
bound}\label{star-capacity-upper-bound}

Fix any d\textgreater=3 arms as in Lemma C, and let b be the projection
of v to I. Let the core C be all legal words with projection either b or
its bitwise complement. For x in C let a(x) count legal edges at x
changing a color in I. Each selected edge of color in I joins exactly
one core vertex to a noncore vertex: it agrees on the other d-1
coordinates, and since d\textgreater=3 the two core projections differ
in more than one coordinate. Each selected edge of color outside I has
both endpoints in C, by Lemma C.

\textbf{Lemma D.} \textbar S\textbar{} \textless= sum\_\{x in C\}
max(1,a(x)).

Proof. Charge each selected I-colored edge to its unique core endpoint.
By Lemma B each selected component is a star. Assign each remaining
selected edge (both endpoints core) to one of its leaf endpoints in that
star, with an arbitrary orientation for a one-edge component. These
recipients are distinct and have no selected I-colored edge, because
their degree in S is one. A core vertex receiving an I-colored edge
receives at most a(x) charges; a recipient of a core-core edge receives
one. QED.

\subsection{4. Exact cyclic counting and finite
reduction}\label{exact-cyclic-counting-and-finite-reduction}

List the chosen coordinates cyclically, at successive distances
g\_0,\ldots,g\_\{d-1\}\textgreater=1, with sum n.~A core projection b is
feasible as a center of a star with those arms exactly when for each
g\_j=1 both adjacent chosen bits b\_j,b\_\{j+1\} are zero (cyclic
indexing). Necessity follows from legality of each flip; sufficiency
follows by setting all unchosen bits zero.

For fixed endpoint bits a,c of a gap g, record p,q, the bits immediately
inside the gap next to the endpoints (for g=1 use p=c,q=a; for g=2 p=q
is the same vertex). Legal assignment counts are:

\begin{itemize}
\tightlist
\item
  g=1: one assignment (p,q)=(c,a) if ac=0; otherwise none.
\item
  g=2: one assignment for each p=q in \{0,1\} with ap=pc=0.
\item
  g\textgreater=3: F\_\{g-1-p-q\} assignments for each p,q in \{0,1\}
  with ap=qc=0. This follows from the adjacency transfer matrix
  {[}{[}1,1{]},{[}1,0{]}{]} raised to power g-2; a zero entry means no
  assignments.
\end{itemize}

Selections in the distinct gaps are independent with the chosen bits
fixed. At chosen coordinate j the flip is legal exactly when
q\_\{j-1\}=p\_j=0. Thus, for a fixed projection b, the total capacity
H(g,b)=sum\_\{x with projection b\} max(1,a(x)) is computed by summing
over all gap records the product of their counts multiplied by max(1,
number of chosen coordinates with both neighbors zero). The total core
capacity is H(g,b)+H(g,1-b).

For any one gap g\_j\textgreater=5, H satisfies the Fibonacci recurrence
as a function of that gap alone: its local factor is F\_\{g\_j-1-p-q\},
and both smaller gaps remain \textgreater=3, so the same p,q records and
factors can be used, with zero counts included. K\_n satisfies the same
recurrence in that gap. Feasibility of b is unchanged when reducing a
gap \textgreater=5 to g\_j-1 or g\_j-2. Therefore margin
K\_n-H(g,b)-H(g,1-b) obeys the same recurrence. Strong induction on sum
g\_j reduces margin nonnegativity to gaps in \{1,2,3,4\}. If one
specified gap is 1, it stays 1 throughout reduction.

\textbf{Theorem E (computer-assisted conditional upper bound).} For
n\textgreater=5, if a selected edge general position set contains a star
with at least four arms, its cardinality is strictly less than K\_n.~If
it contains a three-arm star whose colors include two cyclically
adjacent coordinates, its cardinality is at most K\_n.

Proof. For d=4, the verifier checks all 1831 feasible base tuples, with
nonnegative margin and the sole zero tuple all gaps one and all center
bits zero (n=4). For any reduced tuple with an originally non-base gap,
both smaller sums remain \textgreater=6, so strictness for
n\textgreater=5 follows by induction. For d=3 with one or more cyclic
gaps equal to one, all 64 feasible base tuples have nonnegative margin.
The verifier compares the gap formula against an independent literal
enumeration of every legal word and counts actual legal flips, and
checks center feasibility by actual flips. Lemma D and the recurrence
complete the asserted bounds. QED.

\subsection{5. Finite exact values}\label{finite-exact-values}

The verifier generates forbidden triples in two independent ways: (i)
BFS shortest-path DAGs with propagation of single-edge and two-edge
prefix subsets; (ii) the three possible middle-edge choices and
coordinate monotonicity. Both generate identical forbidden triples. A
triple belongs to a shortest path iff some shortest-path DAG contains
its three edges in order; all possible initial vertices are used. The
monotonicity predicate uses equality of the middle edge and last edge on
the first color, equality of the first and middle edge on the last
color, opposite first/last fixed bits on the middle color, and absence
of repeated coordinate changes outside the three colors. These
conditions suffice because oriented endpoints can be joined by shortest
connectors in the isometric graph.

Given the forbidden triples, the exact search uses an initially complete
pair-compatibility graph on candidate edges. Each included edge makes
two remaining edges incompatible if the three form a forbidden triple.
After including an edge, candidates must also be its neighbors in the
current compatibility graph. This maintains the invariant that a new
candidate has no forbidden triple with two already selected edges, and
future pairs must satisfy all chosen-edge restrictions. A greedy proper
coloring of the current graph gives an upper bound on clique size;
sorting vertices by nondecreasing color and branching in reverse order
gives a bound for the remaining prefix. If selected count plus color
bound or candidate count cannot beat the known witness, the branch is
safely pruned. Otherwise all include/exclude choices are explored
(exclusions proceed by deletion from the prefix). This is an
integer-only exhaustive proof of no larger set, not a numerical solver
assertion.

The computed exact optima are n=3..9: 3,4,7,11,18,29,47, all equal to
K\_n.~The n=5 result is already in the primary reference; no novelty
claimed for it. Exact upper bounds for n\textgreater=10 are not
established here. Verification of witness counts through 20 is only a
consistency check on the analytic lower bound, not an upper-bound proof.

\subsection{6. Remaining work and
literature}\label{remaining-work-and-literature}

For a set using \textless=2 colors, its size is \textless=2F\_\{n-1\},
strictly below K\_n for n\textgreater=5. For a star forest of maximum
degree one, the matching bound
\textbar S\textbar\textless=\textbar V\textbar/2\textless K\_n applies
since
\textbar V\textbar=F\_\{n-1\}+F\_\{n+1\}\textless2(F\_n+F\_\{n-2\}).
Theorem E covers maximum degree \textgreater=4, and a degree-3 star
containing adjacent colors. The remaining cases are star forests of
maximum degree 2 using at least 3 colors, or maximum degree 3 where
every 3-arm star has pairwise nonadjacent colors. Therefore the
conjecture gp\_e(Lambda\_n)=K\_n for all n\textgreater=4 is NOT proved.
n=10 exact search was attempted outside this verifier but did not finish
within the short time limit; it is not an impossibility result. No.75's
frozen manuscript remains unmodified and is not required by the
verification.

Primary reference: S. Klavzar and E. Tan, Edge General Position Sets in
Fibonacci and Lucas Cubes, Bull. Malays. Math. Sci. Soc. 46:120 (2023),
DOI 10.1007/s40840-023-01517-y. Consulted publisher HTML, especially
Proposition 4.1(ii), Theorem 4.6 and Proposition 4.7. Also located
subsequent general-position survey and related papers, with no direct
match for the present general lower construction found in targeted
searches. No exhaustive novelty assessment performed. All mathematical
statements here remain subject to independent review.

Run \texttt{python\ -B\ -I\ verify.py} in Python 3.10+; output should
match EXPECTED\_RESULTS.txt. Standard library only.

\newpage

\section{Frozen v0.2: local star-pair
capacity}\label{frozen-v0.2-local-star-pair-capacity}

\section{No.76 v0.2 --- Local star-pair capacity bounds in Lucas
cubes}\label{no.76-v0.2-local-star-pair-capacity-bounds-in-lucas-cubes}

PRIVATE HOLD / NOVELTY INCONCLUSIVE / REFERENCES PARTIAL. This is an
additive research note, not an all-dimension classification. The exact
frozen v0.1 package is included. No.75 is unchanged. All new general
implications require independent mathematics audit; Windows has not been
run.

\subsection{1. Model and inheritance}\label{model-and-inheritance}

Use the model and Fibonacci indexing in FROZEN\_V01.zip, RESULTS.md.
Write K\_n=F\_n+F\_\{n-2\}. The v0.1 construction gives
gp\_e(Lambda\_n)\textgreater=K\_n for n\textgreater=4, and exact finite
values through n=9 were established there. For n\textgreater=5 the
previous conditional upper bounds cover sets using at most two colors,
matchings, degree \textgreater=4, and three-arm stars with adjacent
colors. We use the self-contained star-forest and projection lemmas B
and C from v0.1, not No.75's main theorem. All following arguments are
for sets S using at least three colors, so the selected graph is a star
forest by Lemma B. A color means a flipped coordinate.

\subsection{2. Two-arm projection and
capacity}\label{two-arm-projection-and-capacity}

The proof of Lemma C in v0.1 works without change for d=2 provided the
edge under consideration has color outside the chosen pair. Indeed,
select the chosen two arms i,j of a star with center v. For an edge f of
color k outside \{i,j\}, a disagreement with v on exactly one of i,j
permits orienting the two-edge path through v to end with the same i,j
bits as f.~Choose the endpoint of f with matching k bit. A shortest
connector changes none of these three coordinates, and the concatenation
changes every coordinate at most once. This is a forbidden geodesic.
Consequently all outside-pair edges lie in the core C consisting of the
two complementary patterns of center bits on the pair. Edges in the pair
have exactly one endpoint in C (one bit flips). The capacity charging
argument of Lemma D therefore works for d=2 also: charge pair-colored
edges to the unique core endpoint, and outside-pair edges to distinct
leaves of their selected star components. A leaf recipient has no
selected pair-colored edge. The bound is the sum over core vertices x of
max(1,a(x)), where a counts actual legal flips in the pair.

The gap counting and recurrence in v0.1 also work for d=2 whenever
n\textgreater=3 (the two cyclic intervals have disjoint interiors). The
new verifier exhausts all feasible gap tuples in \{1,2,3,4\}\^{}2 with
n\textgreater=3 for which the pair is adjacent or the center bits
differ. It compares each fixed-pattern formula against literal
legal-word enumeration and actual flips. All 24 bases satisfy capacity
\textless=K\_n, minimum margin zero. When reducing a gap \textgreater=5,
adjacency status and the center pattern are preserved and both smaller
sums are \textgreater=4. The identical Fibonacci recurrence for target
and capacity proves the conditional upper bound at every
n\textgreater=3. A triple with fewer than three selected colors was
already covered by v0.1 for n\textgreater=5.

\textbf{Theorem Q (conditional).} If S contains a two-arm star with
adjacent colors or unequal center bits on the pair, then
\textbar S\textbar\textless=K\_n.~In particular, to seek a
counterexample at n\textgreater=5 it suffices to consider stars whose
every pair of arms is nonadjacent and has equal center bits. This is a
genuine restriction, not a claim that all remaining configurations
exist.

\subsection{3. Three-arm supplementary
cases}\label{three-arm-supplementary-cases}

Use the original d=3 core capacity of Lemma D. The new verifier tests
all feasible base triples with mixed center projection (not constant) or
with every gap at least 3, and compares against literal enumeration. All
205 bases satisfy the target with minimum margin 4. Gap reduction
preserves these conditions: reducing g\textgreater=5 leaves
g\textgreater=3, and bits remain unchanged. By the recurrence, the
conditional upper bound holds for every such tuple.

\textbf{Theorem R (conditional).} A three-arm star with nonconstant
center projection or with all cyclic gaps at least 3 forces
\textbar S\textbar\textless K\_n.~Together with v0.1, if a
counterexample has maximum degree 3, every three-arm star must have
constant center projection and at least one gap equal to 2; no gap is 1.
This is a necessary condition, not a classification of remaining optima.

\subsection{4. Local capacity by two-arm color
pairs}\label{local-capacity-by-two-arm-color-pairs}

Assume maximum selected degree 2 and that Theorem Q does not apply. Each
two-arm component supplies an unordered pair \{i,j\} with nonadjacent
coordinates and equal center bits. Let P be any nonempty collection of
such pairs, and I the union of their coordinates. The two-arm projection
lemma gives the following constraints simultaneously: - Every selected
edge of color k outside I has both endpoints satisfying equality of bits
on every pair in P. Call this the core condition. - A selected edge of
color k in I satisfies the equality for every pair in P not containing
k. These are only necessary conditions. Ignoring additional forbidden
triples is safe for an upper bound.

Partition legal words into fibers by fixing all coordinates outside I.
Within a fiber, vertices are legal assignments to I. The eligible local
edges are exactly those passing the above necessary conditions. A
selected local subgraph is a forest of stars of maximum degree 2. Let A
be any such local subgraph, e(A) its edge count, and z(A) the number of
core vertices of local degree zero.

\subsubsection{4.1 Arbitrary number of
pairs}\label{arbitrary-number-of-pairs}

Charge every selected edge outside I to a leaf of its selected component
(arbitrary endpoint if a singleton edge). Each recipient is core and has
local degree zero; recipients are distinct. Therefore, adding the local
edge counts and these charges over fibers proves

\begin{verbatim}
2|S| <= sum_fibers max_A [2 e(A)+2 z(A)].
\end{verbatim}

All maxima are finite local optimizations. They are over relaxations,
not assumptions about the existence of globally compatible edges.

\subsubsection{4.2 All two-arm pairs
specified}\label{all-two-arm-pairs-specified}

If P is the complete set of pairs occurring at degree-2 vertices of S,
every exterior edge is isolated in the selected graph: otherwise its
color would belong to a pair in P, contradicting its being outside I.
Count its two endpoints; both are distinct core vertices of local degree
zero. Also require that every local degree-2 vertex uses a pair in P,
since that pair occurs in S. Then

\begin{verbatim}
2|S| <= sum_fibers max_A [2 e(A)+ z(A)],
\end{verbatim}

with this additional local restriction. No exterior edges are assumed to
be actually available at the counted vertices, so the bound remains a
relaxation.

\subsection{5. Cyclic gap summation, recurrence and
certificate}\label{cyclic-gap-summation-recurrence-and-certificate}

Write d=\textbar I\textbar, list I cyclically, and let its cyclic gaps
g\_j be positive with sum n.~For each fixed outside configuration,
record the bits p,q adjacent to the selected endpoints inside each gap.
For g=1 there is no outside vertex; use sentinel p=q=0, and enforce the
adjacency of the two selected bits separately. For g=2 the one outside
bit gives records p=q=0 and p=q=1, each multiplicity one. For
g\textgreater=3 the number of legal internal words with ends p,q is
F\_\{g-1-p-q\}, including zeros. Gaps are independent before assigning
the selected bits. A selected coordinate is blocked from value 1 if one
of its outside neighbors is 1. Gap 1 imposes a direct adjacency between
selected bits. Thus the local graph, core and objective depend only on
the blocked mask and the gap-1 adjacencies, not on other gap lengths.
The total local capacity is obtained by summing over record combinations
with product multiplicities.

For a gap g\textgreater=5, the records can be considered the same four
possible bit pairs, with zero multiplicities allowed, in both smaller
gaps g-1,g-2 (both \textgreater=3). No gap-1 adjacency changes. Hence
the doubled capacity satisfies the Fibonacci recurrence in each such
gap; 2K\_n does too. Restricting pairs to nonadjacent coordinates is
preserved because gaps are only reduced to lengths \textgreater=3. For
n\textgreater=5 and d\textgreater=3, a nonbase gap has smaller total
n-2\textgreater=5; for one nonadjacent pair d=2, the other gap is
\textgreater=2 and n-2\textgreater=5. Strong induction therefore extends
every checked base inequality to all n\textgreater=5 for these cases.

The verifier computes the local optima twice: (i) enumeration of
eligible local edge subsets with exact degree and star checks; (ii)
enumeration of one-edge and two-arm components directly from legal flips
and exact maximum-weight vertex-disjoint packing DP. For (ii), start
from all core vertices credited, then a component on vertex set B with e
edges contributes 2e - factor*\textbar B intersect core\textbar, with
factor 1 or 2 as appropriate. Nonpositive contributions may be omitted
because unused vertices retain their initial credit and no component is
obligatory. Independent gap word enumeration checks the multiplicities
through g=12. Both methods agree on every local optimum used in the
finite cases. All calculations are integers; no numerical solver is
used.

\subsubsection{Theorem S (conditional upper bounds for maximum degree
2)}\label{theorem-s-conditional-upper-bounds-for-maximum-degree-2}

\begin{enumerate}
\def\labelenumi{\arabic{enumi}.}
\tightlist
\item
  If precisely one pair occurs, use section 4.2 with d=2. All 8 feasible
  base cases with n\textgreater=5 have doubled margin \textgreater=2.
\item
  If precisely two distinct pairs occur, either they share a coordinate
  (d=3), or they are disjoint (d=4). Relabel the cyclic order of their
  coordinates; the shared-coordinate case is covered by placing the
  shared coordinate first and using pairs (0,1),(0,2). For disjoint
  pairs there are exactly three pairings of four cyclically ordered
  coordinates, all tested separately. The complete-pair-set bound of
  section 4.2 passes all 36, 144, 255, 144 base cases respectively,
  doubled margins at least 2, 2, 1, 2 respectively.
\item
  If two pairs cross in cyclic order (four distinct coordinates i,j,k,l
  in order, with pairs \{i,k\},\{j,l\}), use section 4.1; no assumption
  on other pairs. All 255 bases have nonnegative doubled margin, minimum
  zero.
\end{enumerate}

Each of these gives \textbar S\textbar\textless=K\_n for all
n\textgreater=5 by the preceding recurrence. Base enumeration excludes
n\textless5 only, since the induction for nonbase tuples never enters
n\textless5. This is a computer-assisted conditional upper theorem, not
a solution of the unrestricted problem.

\subsection{6. Status and limitations}\label{status-and-limitations}

New exact base checks: 1071 total across the eight categories, including
overlap of categories; not 1071 distinct graph instances. Frozen v0.1
contains the earlier 1895 base checks and finite optima through n=9. No
new complete finite optimum is certified in v0.2. The maximum-degree-2
residual case is a family of at least three distinct nonadjacent
equal-center-bit pairs, with no crossings. Maximum degree 3 residual:
every three-arm star has constant center projection, no adjacent arms,
and at least one gap 2; Theorem Q also excludes unequal-bit or adjacent
two-arm stars anywhere in S. These necessary conditions do not imply
existence of counterexamples. This note does not prove the full
conjecture.

No new literature search was performed for the present local-capacity
lemmas. Primary reference and prior searches are documented in frozen
v0.1; novelty remains INCONCLUSIVE. Windows and independent AI/human
audit are pending. Run \texttt{python\ -B\ -I\ verify\_v02.py} and
compare to EXPECTED\_RESULTS.txt; stdlib only. Source package frozen
unchanged; retain the old finite certificates separately. All rights
reserved.

\newpage

\section{Frozen v0.3: completion of the upper
proof}\label{frozen-v0.3-completion-of-the-upper-proof}

\section{No.76 v0.3 --- Exact edge general position number of Lucas
cubes (complete proof
candidate)}\label{no.76-v0.3-exact-edge-general-position-number-of-lucas-cubes-complete-proof-candidate}

PRIVATE HOLD / NOVELTY INCONCLUSIVE / REFERENCES PARTIAL.
Computer-assisted proof candidate, not independently audited or peer
reviewed. No Windows run. No.75 unchanged. The frozen v0.2 archive (and
its nested v0.1) provide the earlier self-contained lemmas, lower
construction, finite exact values, and conditional upper bounds. This
note supplies the two missing upper cases. No claim of extremizer
uniqueness.

\textbf{Theorem (candidate).} For every n\textgreater=3,
gp\_e(Lambda\_n)=K\_n=F\_n+F\_\{n-2\}, with F\_0=0,F\_1=1. Model exactly
as in frozen v0.1. Its Theorem A supplies the lower construction for
n\textgreater=4 and the three-edge star for n=3. Its finite integer
verifier covers n=3..9. It suffices to prove the upper bound for
n\textgreater=5; v0.1 Theorem E and its discussion, and v0.2 Theorems
Q,R,S, are used below.

\subsection{1. Local upper bound for arbitrary maximum selected degree
D=2 or
3}\label{local-upper-bound-for-arbitrary-maximum-selected-degree-d2-or-3}

Consider an edge general position set S using at least three colors. By
frozen Lemma B, S induces a star forest. By Theorem Q, we may assume all
pairs of arms at selected centers are nonadjacent coordinates with equal
center bits; otherwise the desired bound already holds. Let P be a
nonempty subset of the distinct unordered arm pairs that occur anywhere
in S (including the three pairs from a three-arm star), and I their
union. By v0.2 Section 2, each edge of color k must have equal fixed
bits on every pair in P not containing k. In particular, exterior edges
(k outside I) have both endpoints in the core C of words satisfying all
pair equalities. Freeze coordinates outside I to form fibers. Local
eligible edges satisfy the indicated constraints for their color, and
local subgraphs A of S are star forests of maximum degree D. Let e(A)
count their selected edges, z(A) count the legal core vertices not
incident to any local selected edge.

For arbitrary P, charge each exterior selected edge to a leaf endpoint
of its original selected star, with an arbitrary endpoint for a
singleton. These recipients are distinct core vertices, each with local
degree zero. Consequently 2\textbar S\textbar{} \textless= sum\_fibers
max\_A(2e(A)+2z(A)). This upper bound is valid for both D=2 and D=3,
including when exterior edges share a center with local edges.

If P equals the complete set of all occurring arm pairs, every exterior
selected edge is isolated (otherwise its color belongs to P). Its two
distinct endpoints are core vertices of local degree zero. Every local
degree\textgreater=2 vertex has every unordered pair of its incident
colors in P. With this additional local constraint the sharper bound is
2\textbar S\textbar{} \textless= sum\_fibers max\_A(2e(A)+z(A)).

For D=2 and three disjoint pairs P, a simpler upper bound is useful.
Each eligible edge of a color in I changes exactly one of the pairs and
has precisely one core endpoint: bits on the other two pairs agree, and
the changed pair toggles between equal and unequal. Charge each such
selected edge to its unique core endpoint (at most two charges there),
and charge exterior edges as above to distinct core vertices with local
degree zero. Thus 2\textbar S\textbar{} \textless= sum\_fibers
4\textbar C\_fiber\textbar. This avoids all large local packing
computations.

\subsection{2. Recurrence and exact finite
verification}\label{recurrence-and-exact-finite-verification}

List the d=\textbar I\textbar{} coordinates in cyclic order and
g\_j\textgreater=1 their cyclic gaps, sum n.~Use v0.2 Section 5 gap
counts: for g=1 a sentinel 00 with multiplicity 1 and a direct adjacency
constraint; g=2 gives 00,11 each once; g\textgreater=3 gives
F\_\{g-1-p-q\} for p,q in \{0,1\}. The gap records determine the blocked
mask and adjacency constraints on the selected coordinates; local
eligible graphs/core and upper capacities depend only on these data and
P. The total doubled upper bound B(g) is the sum of local values times
product gap multiplicities. For any gap g\_j\textgreater=5 the same four
records can be used at g\_j,g\_j-1,g\_j-2 (zero factors allowed);
adjacency and nonadjacency of pairs are preserved, and
B(g)=B(g-e\_j)+B(g-2e\_j). So M(g)=2K\_n-B(g) satisfies the same
recurrence in each large gap. In our sweeps d\textgreater=3, hence
n-2\textgreater=5 if a gap is \textgreater=5. Induction reduces to gap
tuples in \{1,2,3,4\} with n\textgreater=5 and nonadjacent pairs. This
does not require all legal candidate tuples to occur in actual S.

The verifier independently computes local optimum for d\textless=5 by
two methods. First, enumerate vertex sets of eligible one-, two-, or
three-arm stars and solve exact disjoint vertex packing by subset DP;
begin with all core vertices credited, each component contributes 2
times its edges less the core credit consumed. Omit nonpositive
contributions since no component is mandatory. Second, branch on
eligible edges independently, maintaining selected degrees and neighbor
masks. Reject additions that exceed D, join vertices already of positive
degree, or make a vertex of degree \textgreater=2 adjacent to another
such vertex; these conditions exactly preserve a star forest because the
eligible graph is simple bipartite. When P is complete, reject any new
pair of incident colors outside P. The only bound used in this search is
current incremental score plus twice the remaining edge count; selecting
another edge can increase score by at most 2. Both methods agree at
every local instance used. For d=6 (three disjoint pairs) use the
preceding simple core bound, directly checking the unique-core-endpoint
property. Gap masks are computed both from Fibonacci counts and by
enumerating literal short gap words with OR convolution. No floating
point or external solver.

\subsection{3. Completion of maximum degree
two}\label{completion-of-maximum-degree-two}

If there are at most two distinct arm pairs, or there is a crossing pair
of pairs, frozen Theorem S applies. Otherwise choose any three distinct
occurring pairs. They are noncrossing with four distinct alternating
endpoints excluded, and all pairs nonadjacent and equal-bit at their
centers. Their union has d between 3 and 6. Enumerate all three-element
subsets of the unordered pairs of d labeled positions whose union is all
d, reject crossing pairs, and quotient by cyclic rotations and
reflection (preserving cyclic order or reversing it). This produces
exactly eight representative patterns, listed by verify\_v03.py; the
rotation/reflection also permutes the gaps and legal assignments, so the
canonical pattern suffices. Apply Section 1 arbitrary-P bound with D=2,
using the simpler core bound at d=6. All 5,463 bases have
M\textgreater=0. Section 2 proves the upper bound for every nonbase
tuple as well.

\subsection{4. Completion of maximum degree
three}\label{completion-of-maximum-degree-three}

Select a three-arm star and let T be its three coordinates. By Theorem Q
every pair in T has equal center bits and is nonadjacent. P\_0 consists
of the three pairs of T. If no other arm pair occurs anywhere in S, use
Section 1 complete-P bound with D=3 and P=P\_0. All 27 bases pass.

Otherwise choose any occurring pair p not in P\_0 and use P=P\_0 union
\{p\}. Its union has d=4 or 5, since p has at least one coordinate
outside T. Unlike Section 3 there is NO noncrossing restriction.
Enumerate every choice of T and p whose union is all d, and reduce under
dihedral relabeling. Exactly four patterns remain (two for d=4, two for
d=5). Use arbitrary-P bound with D=3. Among 1,452 bases exactly one has
negative margin, -2: d=5, P=((0,1),(0,3),(1,3),(2,4)), g=(2,1,1,1,1),
n=6. This is a failure of an upper-bound relaxation, NOT a
counterexample to the theorem. The frozen v0.1 exact finite certificate
for n=6 gives gp\_e=11=K\_6. Crucially, the exceptional tuple has every
gap \textless=2. The recurrence is only used at gaps \textgreater=5 and
its children have that gap \textgreater=3. Therefore no recurrence step
can have this exceptional tuple as a child, and no sequence of
recurrence steps starting at a nonbase tuple can terminate at it: at
least one initially large gap stays \textgreater=3. All other bases have
nonnegative margins, so strong induction proves nonnegativity at every
nonbase tuple. The exceptional finite case is covered directly by the
frozen n=6 certificate. Total D=3 bases 1,479, including the exceptional
base; 1,478 nonnegative. This completes D=3.

\subsection{5. Conclusion and validation
scope}\label{conclusion-and-validation-scope}

For n\textgreater=5 sets with \textless=2 colors are bounded by
2F\_\{n-1\}\textless K\_n; with at least three colors Lemma B gives a
star forest. Maximum degree 1 is bounded by
\textbar V\textbar/2\textless K\_n; \textgreater=4 is covered by frozen
Theorem E; degrees 2 and 3 are covered above. The empty set is trivial.
Small n=3,4 have frozen exact certificates. Together with Theorem A the
asserted equality follows, conditional on correctness of these arguments
and verified finite bases. The identity K\_n=2F\_\{n-1\}+F\_\{n-4\} for
n\textgreater=4 yields the strict lower-bound improvement for
n\textgreater=5. The 2023 primary paper establishes n=5 already, so that
finite value is not claimed new. This is a candidate classification of
Lucas cubes only, not a dependence on or alteration of No.75's main
theorem.

New finite sweep 6,942 cases (not all distinct actual graphs), one
negative relaxation base explicitly handled by exact finite certificate.
New verifier checks only Section 2--4 bases. Run frozen v0.1 verify.py
and v0.2 verify\_v02.py as well to reproduce all dependencies; compare
each to its own EXPECTED\_RESULTS.txt. All packages use standard library
only, Python 3.10+.

Primary prior work: S. Klavzar and E. Tan, Edge General Position Sets in
Fibonacci and Lucas Cubes, Bull. Malays. Math. Sci. Soc. 46:120 (2023),
DOI 10.1007/s40840-023-01517-y. Prior work reports the 2F\_\{n-1\} lower
bound and dimension-five exception. This note does not assert that the
new equality was an explicit conjecture of that paper. No fresh
literature search was performed in v0.3; full novelty is INCONCLUSIVE.
Independent audit and Windows reproduction required before publication.
All rights reserved.

\newpage

\section{v0.4: proof-dependency
clarification}\label{v0.4-proof-dependency-clarification}

\section{No.76 v0.4 --- Proof-dependency clarification (addendum to
frozen
v0.3)}\label{no.76-v0.4-proof-dependency-clarification-addendum-to-frozen-v0.3}

PRIVATE HOLD / NOVELTY INCONCLUSIVE / REFERENCES PARTIAL / NOT PEER
REVIEWED / All rights reserved.

The sole mathematical parent is FROZEN\_V03.zip, including its nested
v0.2 and v0.1. No parent text, code, theorem, scope, license or result
is changed. The additional exact finite check is corroboration, not a
replacement for the frozen proofs. No claim of uniqueness is added. This
note clarifies the proof dependencies in response to the independent AI
audit report on 2026-11-19.

\subsection{Statement (unchanged)}\label{statement-unchanged}

For every integer n \textgreater= 3, gp\_e(Lambda\_n) = K\_n = F\_n +
F\_\{n-2\}, where F\_0=0 and F\_1=1. The proof is computer-assisted and
remains a candidate pending confirmation of this clarification by the
independent auditor.

\subsection{Exact dimension dependencies vs capacity
bases}\label{exact-dimension-dependencies-vs-capacity-bases}

The whole-graph exact optima in frozen v0.1 Section 5 and the finite
capacity inequalities in the three frozen versions are different types
of evidence. The following is a dependency audit, not a removal of any
required capacity check.

\begin{enumerate}
\def\labelenumi{\arabic{enumi}.}
\tightlist
\item
  Lower bound: v0.1 Theorem A (including its n=3 observation) supplies
  the construction in all dimensions.
\item
  n=3 and n=4: v0.3 Section 5 explicitly invokes the frozen exact
  certificates. These two whole-graph optima are required by the proof
  as written.
\item
  n\textgreater=5: a set using \textless=2 colors is bounded by v0.1
  Section 6; otherwise Lemma B yields a star forest. The matching bound
  handles maximum selected degree 1. Theorem E (v0.1 Section 4) handles
  maximum degree \textgreater=4. For degrees 2 and 3, Theorem Q (v0.2
  Section 2) handles adjacent or unequal-bit pairs. The remaining
  degree-2 cases are covered by Theorem S (v0.2 Section 5) and v0.3
  Section 3. The remaining degree-3 cases are covered by v0.3 Section 4.
\item
  The sole negative relaxation margin in v0.3 Section 4 is -2, at d=5,
  P=((0,1),(0,3),(1,3),(2,4)), gaps=(2,1,1,1,1), n=6, up to symmetry.
  This case needs the exact whole-graph bound gp\_e(Lambda\_6)=11.
  Reducing any gap \textgreater=5 leaves it \textgreater=3 in both
  children, so neither child equals the exceptional all-gaps\textless=2
  tuple. Induction on the sum of gaps for nonexceptional tuples is
  valid: for every nonbase parent, both children are nonexceptional and
  smaller. The exceptional tuple is treated separately by the n=6 bound;
  its negative relaxed margin is never used to infer a nonnegative
  margin.
\item
  No other whole-graph exact optimum is invoked by this argument. In
  particular, the n=8 and n=9 whole-graph optima are corroborating
  computations in frozen v0.1 Section 5, not premises of the general
  equality. Neither are the n=5 and n=7 exact optima required (n=5 is
  independently known in the cited primary paper). This does NOT waive
  capacity bases with total dimension 8 or 9 or any other dimension.
  Preserve all 1,895 + 1,071 + 6,942 = 9,908 capacity checks (including
  overlaps) and their recurrence and local-optimization proofs. Theorem
  R is an auxiliary, stronger conditional restriction, not needed for
  the final case split as written; its certificate is still preserved.
\end{enumerate}

\subsection{Additional integer-only
corroboration}\label{additional-integer-only-corroboration}

check\_required\_dimensions.py is a new independent-style implementation
by the development side. It enumerates legal words and edges, calculates
all-pairs graph distances via BFS, and tests every edge triple using all
six orders and eight orientations. An ordered oriented triple lies on a
geodesic iff the endpoint distance equals 3 plus the two connector
distances: necessity follows from geodesic subpaths; sufficiency follows
by concatenation of shortest connectors. The feasible witness is checked
directly against these triples. The search then enumerates increasing
edge subsets; each inclusion disallows any later edge forming a
forbidden triple with two chosen edges; only remaining-cardinality
pruning is used. This is a complete exact search for a set of target+1
edges, which suffices by heredity. No external solver, coloring bound,
frozen triplet generator or author-auditor implementation is imported.
It confirms n=3,4,6 and the nonexistence of a 12-edge set for n=6. It
does not rerun independent n=8,9 optimization or repeat the entire
independent audit.

\subsection{Status}\label{status}

The supplied independent report passed the model, lower construction,
structural lemmas, induction, and separately established exact optima
for n=3..7. Its final theorem verdict was conditional, with n=8,9
independent optimum proofs unfinished; its publication recommendation
was HOLD. Those verdicts are preserved pending an explicit delta audit.
The frozen original verifiers, including their n=8,9 searches, were all
replayed on Linux with byte-identical expected outputs during this
addendum. Earlier Windows FULL PASS pertains to v0.1--v0.3, not the new
addendum checker. No fresh Windows execution or fresh literature
assessment was performed in v0.4. No claim of direct novelty, formal
verification or human peer review is made.

\end{document}
